\clearpage\SourceLeaf{87}%
\pdfbookmark[1]{Noch etwas über die Bestimmung des Osterfestes}{osterfest1807}
\begin{center}
{\large NOCH ETWAS ÜBER DIE BESTIMMUNG DES OSTERFESTES.\par}
\medskip\hrule\medskip
{\small Braunschweigisches Magazin. 1807. September 12.\par}
\medskip\hrule
\end{center}\medskip
Die zyklische Berechnung des Osterfestes, von der in diesen Blättern schon mehrmals die Rede war, ist zwar an sich leicht, beruht aber auf Grundlagen, die nicht jedem bekannt sind. Sie setzt außerdem Hilfstafeln voraus, die nicht jeder zur Hand hat und die, wie das Beispiel des Herrn Superintendenten \textsc{Helmuth} gezeigt hat, selbst in gedruckten Schriften fehlerhaft sein können. Trotzdem möchte man nicht selten das Osterdatum für ein künftiges Jahr oder für ein vergangenes Jahr wissen, für das man keinen Kalender zur Hand hat. Manchem mag daher ein besonderes Verfahren zur Bestimmung des Osterfestes willkommen sein, bei dem man die Grundlagen und Fachwörter der zyklischen Berechnungsart überhaupt nicht kennen muss, keine Hilfstafeln braucht und nur ein paar ganz einfache Rechenoperationen auszuführen hat. Ich habe diese Methode bereits vor sieben Jahren in der von Herrn \textsc{von Zach} herausgegebenen \textit{Monatlichen Correspondenz} (Augustheft 1800) mitgeteilt. Da sie aber wohl nur wenigen Lesern dieses Magazins bekannt geworden ist, mag es dem einen oder anderen helfen, hier das Wesentliche dieses Verfahrens noch einmal erklärt zu finden. Ich beschränke mich hier jedoch darauf, die Regeln zur Bestimmung des Osterfestes im gregorianischen Kalender für das gegenwärtige Jahrhundert, die unmittelbar %
\SourceLeaf{88}%
folgenden und die vergangenen so darzustellen, dass selbst Ungeübte sie verstehen können. Dazu werde ich noch die Regeln für dieselbe Rechnung im julianischen Kalender angeben. Die Erweiterung der Regeln auf die späteren Jahrhunderte und ihren Zusammenhang mit der ursprünglichen Einrichtung unseres Kalenders muss man an der angegebenen Stelle nachlesen. [S.~75 d.~B.]

Zur Bestimmung des Osterfestes für jedes beliebige Jahr des gegenwärtigen Jahrhunderts von 1800 bis einschließlich 1899 sind die folgenden Regeln zu beachten.

$1^{0}$. Man teilt die Jahreszahl durch 19. Der Quotient bleibt ganz unberücksichtigt. Man merkt sich nur den Rest. Um ihn von den Resten zu unterscheiden, die bei den folgenden Operationen entstehen, wollen wir ihn den \textit{ersten} Rest nennen. Geht die Division auf, nimmt man 0 als Rest. Das gilt auch bei den übrigen Divisionen.

$2^{0}$. Man multipliziert diesen ersten Rest mit 19, addiert 23 zum Produkt und teilt die Summe durch 30. Auch hier merkt man sich nur den Rest, den wir den \textit{zweiten} Rest nennen.

$3^{0}$. Man teilt die Jahreszahl durch 4. Der Rest heißt der \textit{dritte} Rest.

$4^{0}$. Man bildet die Summe aus der vierfachen Jahreszahl, dem sechsfachen zweiten Rest, dem doppelten dritten und der Zahl 4. Bei der Division dieser Summe durch 7 erhalten wir den \textit{vierten} Rest.

$5^{0}$. Die Summe des zweiten und vierten Restes gibt nun an, wie viele Tage nach dem 22. März Ostern eintritt. Ist diese Summe höchstens 9, so fällt Ostern noch in den März. Das Datum erhält man, indem man zu dieser Summe 22 addiert. Ist die Summe größer als 9, so fällt Ostern in den April. Das Datum erhält man, indem man von der Summe 9 abzieht.

Ich will diese Regeln an ein paar Beispielen erläutern. Zuerst sei Ostern für das Jahr 1801 zu berechnen. Hier ergibt

\textit{nach der ersten Regel} die Division durch 19 den Quotienten 94, um den wir uns nicht kümmern, und den Rest 15, der also der erste Rest ist.

\textit{Nach der zweiten Regel} ist das Neunzehnfache des ersten Restes 285. Addiert man 23, ergibt sich 308. Bei der Division durch 30 bleibt 8 als zweiter Rest.

\textit{Nach der dritten Regel} ist der dritte Rest 1.
\SourceLeaf{89}%

\begin{center}\small
\begin{tabular}{llr}
\textit{Nach der vierten Regel} & addieren wir die vierfache Jahreszahl &7204\\
& den sechsfachen zweiten Rest &48\\
& den doppelten dritten Rest\quad . &2\\
& und die Zahl\quad .\ .\ .\ .\ .\ . &4\\\cline{3-3}
&\hfill ergibt die Summe &7258
\end{tabular}
\end{center}
Teilt man die Summe durch 7, bleibt 6 als vierter Rest.

\textit{Nach der fünften Regel} erhalten wir schließlich 14 als Summe des zweiten und vierten Restes. Da sie um 5 größer ist als 9, fällt Ostern auf den 5. April. — Der Gründonnerstag dieses Jahres, an den die Ereignisse des Tages jetzt wieder erinnern, fiel also auf den 2. April.

Für das Jahr 1808, das Herrn Superintendenten \textsc{Helmuth} den Anlass zu seinem Aufsatz gab, ergibt sich auf dieselbe Weise Folgendes.

Nach der ersten Regel ist der Rest 3.

Nach der zweiten Regel ist das Neunzehnfache des ersten Restes 57. Addiert man 23, erhält man 80, also ist der zweite Rest 20.

Nach der dritten Regel ist der dritte Rest 0.

Nach der vierten Regel haben wir 7232, 120, 0 und 4 zu addieren. Das ergibt die Summe 7356, also den vierten Rest 6.

Nach der fünften Regel erhalten wir die Summe 26. Zieht man davon 9 ab, ergibt sich der 17. April als Osterdatum.

Nach diesen Erläuterungen kann jeder, der die vier Grundrechenarten kennt, die Rechnung für andere Jahre ohne Weiteres ausführen.

In der oben angegebenen Form gelten diese Regeln für die hundert Jahre von 1800 bis 1899. Für andere Jahrhunderte ändern sie sich etwas. Im vergangenen Jahrhundert von 1700 bis 1799 besteht der einzige Unterschied darin, dass man in der Summe nach der vierten Regel die Zahl 4 durch 3 ersetzt. Im siebzehnten Jahrhundert dagegen, oder genauer von 1583 (als Ostern erstmals nach der neuen Einrichtung bestimmt wurde) bis 1699, ist in der eben genannten vierten Regel die Zahl 4 durch 2 zu ersetzen. Außerdem wird in der zweiten Regel die Zahl, die zum 19-fachen ersten Rest addiert wird, von 23 auf 22 geändert. Alles Übrige bleibt gleich. Nur diese beiden Punkte betreffen auch die Änderungen für die unmittelbar folgenden Jahrhunderte. Anstelle der Zahl 23, die in der %
\SourceLeaf{90}%
ersten Regel für das achtzehnte und neunzehnte Jahrhundert gilt, setzt man im zwanzigsten, einundzwanzigsten und zweiundzwanzigsten die Zahl 24, im dreiundzwanzigsten 25, im vierundzwanzigsten 26 und im fünfundzwanzigsten wieder 25. Anstelle der Zahl 4, die in der vierten Regel im neunzehnten Jahrhundert gebraucht wird, oder der 3, die im achtzehnten gebraucht wird, nimmt man
\begin{center}\small
\begin{tabular}{lr}
im zwanzigsten und einundzwanzigsten\quad .\ .\ .\ . &5\\
im zweiundzwanzigsten\quad .\ .\ .\ .\ .\ .\ .\ .\ . &6\\
im dreiundzwanzigsten\quad .\ .\ .\ .\ .\ .\ .\ .\ .\ . &0\\
im vierundzwanzigsten und fünfundzwanzigsten &1
\end{tabular}
\end{center}
Die allgemeinen Regeln für diese Änderungen findet man an der angegebenen Stelle der Mon. Corresp. [S.~78 d.~B.]

Ich habe nur noch zu bemerken, dass die Epaktentafel im gregorianischen Kalender so eingerichtet ist, dass es in einzelnen Jahren zwei Arten von Ausnahmen von den obigen Regeln geben kann. Die erste besteht darin, dass jedes Mal, wenn der zweite Rest 29 und der vierte Rest 6 ist und Ostern nach der fünften Regel also auf den 26. April fallen müsste, das Fest eine Woche früher gefeiert wird, also am 19. April. Dieser Fall ist bisher nur einmal eingetreten, nämlich im Jahr 1609, und wird erst im Jahr 1981 wieder eintreten. Die zweite Ausnahme kann eintreten, wenn der zweite Rest 28 und der vierte 6 ist. Nach der fünften Regel müsste Ostern dann am $25^{\mathrm{sten}}$ April sein, wird aber tatsächlich auf den 18. April verlegt, und zwar nur, wenn der erste Rest nicht unter 11 lag. In diesem Fall tritt die Ausnahme nicht ein, und Ostern bleibt am 25. April. Diese zweite Ausnahme ist bisher noch nie eingetreten und wird zum ersten Mal erst im Jahr 1954 vorkommen.

Alles Vorangehende betrifft die Berechnung des Osterfestes im gregorianischen Kalender. Für den julianischen Kalender brauchen die obigen fünf Regeln nur in zwei Punkten geändert zu werden. In der zweiten Regel ersetzt man die Zahl 23 durch 15 und in der vierten Regel 4 durch 6. Mit diesen Änderungen gelten die Regeln allgemein für alle Jahrhunderte. Auch diese Rechnung will ich an einem Beispiel erläutern. Für das Jahr 1808 erhielten wir nach der ersten Regel, die unverändert bleibt, den ersten Rest 3. Multipliziert man ihn mit 19 und addiert 15, ergibt sich 72. Bei der Division durch 30 bleibt %
\SourceLeaf{91}%
12 als zweiter Rest. Der dritte Rest ist wie oben 0. Nach der vierten Regel haben wir nun 7232, 72, 0 und 6 zu addieren. Das ergibt 7310, also den vierten Rest 2. Seine Summe mit dem zweiten Rest ist 14. Zieht man nach der fünften Regel 9 davon ab, bleibt 5. Ostern fällt also auf den 5. April nach dem julianischen Stil, der dem 17. April nach dem gregorianischen entspricht. In diesem Jahr fällt Ostern also zufällig in beiden Kalendern auf denselben Tag, obwohl das julianische Datum um 12 Tage niedriger ist.

Übrigens lässt sich zeigen, dass im 143. Jahrhundert, das heißt von 14200 bis 14299, Ostern in beiden Kalendern immer auf dasselbe Datum fallen würde, wenn beide Kalender so lange unverändert in Gebrauch blieben, was freilich kaum zu erwarten ist. Dasselbe Datum läge dann aber in den beiden Kalendern immer 15 volle Wochen auseinander, und das julianische Osterfest würde also 15 Wochen später eintreten als das gregorianische. Mit dem Jahr 14300 endet diese Übereinstimmung wieder. Sie beginnt 43200 von Neuem. Dann würde der Unterschied aber 46 Wochen betragen.

Schließlich bemerke ich, dass im Maiheft der Mon. Corresp. von 1802 [S.~80 d.~B.] auch eine rein arithmetische Regel zur Berechnung des Pessachfestes angegeben ist, für deren Anwendung keine weitere Kenntnis der Einrichtung des jüdischen Kalenders nötig ist.

\begin{center}\rule{.25\linewidth}{.3pt}\end{center}
